\documentclass[12pt]{article}
\usepackage{styles/cwilliams-standard}

% Set up document metadata
\setauthors{Tzu-Chen Chiu, Arjun Sivakumar, Mara Schimmel, Christopher Williams}
\setheaderauthors{Chiu, Sivakumar, Schimmel, Williams}
\settitle{Sprint 1 Report - Semester 2}
\setclass{IPD}

\begin{document}

\maketitlepage

\section{Original Sprint 1 Backlog}

\begin{figure}[H]
    \centering
    \includegraphics[width=1\textwidth]{images/sprint_4_backlog.png}
    \caption{Spring 4 Backlog}
\end{figure}

\section{Sprint 4 Backlog - End of Sprint Status}

\begin{figure}[H]
    \centering
    \includegraphics[width=1\textwidth]{images/sprint4_backlog_end.png}
    \caption{Spring 4 Backlog - Updated}
    \label{fig:sprint4updated}
\end{figure}

\section{Project Management Screenshots}

The Sprint backlogs in Google Sheets are being used to manage tasks.

\section{Sprint Retrospective}

\subsection*{Mara}

For this sprint, my goal was to gain a better understanding of transformers and how they work. 
In addition, I got an additional task to visualize the transformer results using Streamlit, 
a Python library, and began working with small data sets to choose what version we prefer. 
For the next sprint, I plan on learning more about transformers, setting up the ARC for 
myself, and continuing to work on the data visualization.

\subsection*{Chris}

During this sprint, I focused on fixing the embedding pipeline according to the assigned 
tasks and collaborated closely with Tzu Chen to improve the model’s performance. A significant 
portion of the sprint was spent analyzing Google’s first and second versions of the 
transformer-based quantum decoder (AlphaQubit) to better understand architectural changes 
and design choices that could inform our implementation. One major impediment was the 
reliance on a personal computer for running experiments, which made training and evaluation 
time-consuming and limited iteration speed. Despite this constraint, collaboration went well 
through frequent discussions and shared analysis, and the sprint resulted in clearer 
architectural direction. Moving forward, earlier access to adequate compute resources 
would significantly improve productivity and experimentation speed.

\subsection*{Tzu-Chen}

We have built the basic Transformer based on Google’s pseudo code. We studied which features 
should be input to the model: stabilizer id, syndrome(hard input in pretraining), events and 
so on. Then embedding all features into the matrix. We then turn the matrix into four multihead 
attention matrices, which is the input being fed to the Transformer model later. We 
recorded the state of the model result by RNN. After all cycles are finished running, 
we could calculate LER and compare with the baseline MWPM. Overall, it is a very huge 
task to complete. Thanks to my teammates and mentor’s assistance, the model has the 
initial architecture now. For the next sprint, we plan to mimic the dataset from Google, 
then we can actually start training our model.

\subsection*{Arjun}

During this sprint, my original task was to integrate an adapter to work with real-time data, 
however the project scope shifted and real-time data was no longer part of our workflow. As a 
result, I focused on enabling our models to run on Virginia Tech’s ARC compute cluster, which 
involved obtaining the correct allocation, writing and debugging SLURM job scripts, configuring 
the environment, and ensuring results could be generated and stored reliably. I encountered 
technical impediments related to allocation access, which required coordinating with and 
emailing the WISE team to resolve. Beyond technical work, I supported my teammates by 
actively participating in discussions, sharing progress updates, and helping troubleshoot 
issues where possible. Overall, the sprint went well in terms of infrastructure setup 
and collaboration, though earlier clarification of changing project requirements could 
have reduced rework and improved efficiency.

\section{Product Backlog (Updated)}

\begin{figure}[H]
    \centering
    \includegraphics[width=1\textwidth]{images/sprint2_product_backlog_sem2.png}
    \caption{Updated Product Backlog}
\end{figure}

\section{Issue Tracking}

\begin{table}[H]
\centering
\begin{tabular}{@{}llll@{}}
\toprule
\textbf{Issue} & \textbf{Description} & \textbf{Status} & \textbf{Date Completed} \\ \midrule
ARC Access & Getting all members access to ARC & Ongoing & Ongoing \\
Ensuring data quality & Ensure correct data & Ongoing & Ongoing \\
Creating better architecture & $\leftarrow$ & Complete & 2/5/2026 \\
\end{tabular}
\caption{Issue Tracking Summary}
\end{table}

\section{Sprint 2 Backlog - Semester 2}

\begin{figure}[H]
    \centering
    \includegraphics[width=1\textwidth]{images/sprint2_backlog_sem2.png}
    \caption{Spring 4 Backlog}
\end{figure}

\end{document}